% Hardware, software, and the plan of measurements of the Day 9 lab.
\begin{frame}{Hardware and software}
  \begin{columns}[T]
    \begin{column}{0.46\textwidth}
      \small
      \textbf{Hardware for each group}
      \begin{itemize}
        \item XIAOML Kit and a USB-C data cable. The lab uses no sensor and
              no display.
        \item Raspberry Pi 5 (8 GB) with the active cooler, the card of your
              group, and the 27 W power supply.
        \item Laptop in the Wi-Fi \code{edgeai-lab}.
        \item A USB power meter, if the classroom has it.
      \end{itemize}
    \end{column}
    \begin{column}{0.50\textwidth}
      \small
      \textbf{Software}
      \begin{itemize}
        \item Arduino IDE with the esp32 core 3.3.12 and the library
              Chirale\_TensorFlowLite 2.0.0.
        \item On the board: \code{\textasciitilde/tflite\_env} (LiteRT, ONNX
              Runtime) and \code{\textasciitilde/yolo} (LiteRT, NCNN).
        \item The model files of Days 7 and 8 on the board.
        \item Python 3 on the laptop for \code{pi/make\_report.py}. It needs
              no package.
      \end{itemize}
    \end{column}
  \end{columns}
  \vskip 0.4em
  \small
  \renewcommand{\arraystretch}{1.15}
  \begin{tabular}{@{}L{0.24\textwidth}L{0.70\textwidth}@{}}
    Lab files & \code{Labs/day09/} on the laptop,
                \code{\textasciitilde/edgeai/day09/} on the board \\
    The models & The four reference models of MLPerf Tiny (Part 1 of the
                 lecture): keyword spotting and image classification \\
  \end{tabular}
\end{frame}

\begin{frame}{The same model file on two boards}
  \begingroup\centering
  \begin{tikzpicture}[font=\footnotesize, >={Stealth[length=5pt]},
      b/.style={draw=coolgray, rounded corners=3pt, fill=boxgray,
                minimum width=2.0cm, minimum height=1.0cm, align=center},
      lab/.style={font=\scriptsize, text=coolgray, align=center}]
    \node[b, draw=darkcyan] (a) at (0,0) {Protocol\\(Part A)};
    \node[b, draw=treegreen] (b) at (2.9,0.85) {XIAOML Kit\\(Part B)};
    \node[b, draw=darkorange] (c) at (2.9,-0.85) {Raspberry Pi\\(Part C)};
    \node[b, draw=cardinalred] (d) at (5.8,0) {Tables\\(Part D)};
    \node[b] (e) at (8.6,0) {Report};
    \draw[->] (a) -- (b); \draw[->] (a) -- (c);
    \draw[->] (b) -- (d); \draw[->] (c) -- (d); \draw[->] (d) -- (e);
    \node[lab, anchor=south] at (2.9,1.4) {4 models, \code{int8} and
      \code{float32}};
    \node[lab, anchor=north] at (2.9,-1.4) {the same 4 models, MobileNetV2,
      your detector};
    \node[lab, anchor=north] at (5.8,-0.55) {energy,\\battery};
  \end{tikzpicture}\par\endgroup
  \vskip 0.4em
  \small
  \begin{itemize}
    \item The kit runs the four files with TensorFlow Lite Micro. The
          Raspberry Pi runs the same four files with LiteRT. The two boards
          get the same test input.
    \item The Raspberry Pi also measures the models of Days 7 and 8: two
          runtimes and two thread counts for each model.
    \item The window is the model only on the two boards.
  \end{itemize}
  \keypoint{A comparison of two boards is valid only with the same model
    file, the same input, and the same window.}
\end{frame}
